\documentclass[12pt]{article}
\usepackage[utf8]{inputenc}
\usepackage[T1]{fontenc}
\usepackage{amsmath,amssymb}
\usepackage{geometry}
\geometry{margin=2.5cm}

\begin{document}

\noindent\underline{\textbf{Subgrup normal.\ Grup factor.}}

\bigskip
\noindent Câte transv.\ la stânga ale lui $H$ în $S_n$ există? \ $=\big((n-1)!\big)^n$.\footnote{cuvântul de aici (,,det.'') e neclar în original.}

\bigskip
\noindent\underline{Lema 1.} Fie $G$=grup, $N\leq G$.

\noindent\underline{V.A.S.E}: [Vom arăta sunt echivalente:]

\begin{enumerate}
\item[(1)] $(\forall)\, a\in G, (\forall)\, x\in N \Rightarrow axa^{-1}\in N$.
\item[(2)] $(\forall)\, a\in G \Rightarrow aN=Na$
\item[(3)] $R_N^s = R_N^d$ ($\overset{\text{not}}{=} R_N$).
\end{enumerate}

\noindent Un subgrup cu aceşti proprietăți va fi numit \underline{subgrup normal} al lui $G$

\noindent (\underline{divizor normal}).

\bigskip
\noindent\underline{Dem}: $\big[G$=comutativ $\Rightarrow (\forall)$ subgrup e normal$\big]$

\bigskip
\begin{center}
(2) $\Rightarrow$ (1) $\Rightarrow$ (3) $\Rightarrow$ (2)\footnote{în original acest ciclu de implicații e desenat ca o diagramă triunghiulară cu săgeți, nu liniar.}
\end{center}

\bigskip
\noindent\underline{(2) $\Rightarrow$ (1)}, \ $x\in N$.

\[
ax\in aN=Na
\]
\[
\Rightarrow \exists\, x_1\in N \text{ a.î. } ax=x_1a \ \Rightarrow \ axa^{-1}=x_1
\]
\[
\Rightarrow axa^{-1}\in N.
\]

\bigskip
\noindent\underline{(3) $\Rightarrow$ (2)}

\[
aN = [a]_{R_N^s} = [a]_{R_N^d} = Na
\]

\bigskip
\noindent\underline{(1) $\Rightarrow$ (3)}

\noindent\footnote{fragmentul introductiv de aici e neclar în original.} $(x,y)\in R_N^s \Leftrightarrow x^{-1}y\in N \Leftrightarrow y^{-1}x\in N$

\[
= x(y^{-1}x)x^{-1} \overset{(1)}{\Longrightarrow} xy^{-1}\in N \Rightarrow (x,y)\in R_N^d
\]

\bigskip
\noindent\underline{Exp}: (normal) Grup factor.

\noindent (1) $(\mathbb{Z},+,0)$, $H=n\mathbb{Z}$, $n\geq 0$.

\noindent $H$ normal în $\mathbb{Z}$, \ $H\trianglelefteq\mathbb{Z}$.

\noindent (2) $(S_3,\circ)$. \quad $\sigma = \begin{pmatrix}1&2&3\\2&3&1\end{pmatrix}$ \quad $\sigma^2=\begin{pmatrix}1&2&3\\3&1&2\end{pmatrix}$, \ $\sigma^3=\begin{pmatrix}1&2&3\\1&2&3\end{pmatrix}$.

\end{document}
